% TOP_PROBLEM: {"id": 30006963, "title": "Zariski Dense Orbit Conjecture", "category_id": 11, "category": {"id": 11, "name": "dynamical_systems", "display_name": "Dynamical Systems", "description": "Problems about long-term behavior of deterministic systems, Hamiltonian dynamics, and periodic orbits.", "slug": "dynamical-systems", "order_index": 11, "created_at": "2026-07-31T15:26:25.671Z"}, "status": "open"}
% FIRST_POSTED: 2026-09-27
% !TeX program = lualatex
% ATTEMPT_STATUS: unresolved
% Model/process: GPT_6_Astra_Ultra; detailed mathematical attempt,
% primary-source check, exact arithmetic diagnostics, adversarial review.
% Prepared 27 September 2026; no general resolution or novelty claimed.
\documentclass[11pt]{article}
\usepackage[margin=25mm]{geometry}
\usepackage{fontspec}
\setmainfont{Latin Modern Roman}
\usepackage{amsmath,amssymb,mathtools}
\usepackage{unicode-math}
\setmathfont{Latin Modern Math}
\usepackage{xurl,hyperref}
\hypersetup{colorlinks=true,urlcolor=blue}
\setcounter{secnumdepth}{0}
\setlength{\parindent}{0pt}
\setlength{\parskip}{5pt}
\setlength{\emergencystretch}{3em}
\begin{document}
\section*{\texorpdfstring{Zariski Dense Orbit Conjecture}{Zariski Dense Orbit Conjecture}}\label{30006963}
\subsection{Definitions and mathematical statement}
Let $X$ be an irreducible quasiprojective variety over an algebraically closed characteristic-zero field $K$, and let $f:X\dashrightarrow X$ be dominant. The invariant rational functions form
\[
 K(X)^f=\{g\in K(X):g\circ f=g\}.
\]
Assume $K(X)^f=K$. The conjecture asks for a point $x\in X(K)$ whose entire forward orbit is defined and Zariski dense:
\[
 f^n(x)\notin\operatorname{Indet}(f)\ (n\ge0),
 \qquad\overline{\{f^n(x):n\ge0\}}^{\rm Zar}=X.
\]
The hypothesis rules out invariant rational fibrations detected by a nonconstant rational function. The field need not be uncountable, so an argument simply avoiding a countable union of proper closed sets is not sufficient in the full scope. The conclusion is existence of one point, not a dense set of such points.
\par\smallskip\textit{Formulation and concept references:} \href{https://personal.math.ubc.ca/~dghioca/papers/conjecture_automorphisms.pdf}{[S1]}.

\subsection{Short English statement}
If a dominant rational self-map has no nonconstant invariant rational function, must it have a well-defined Zariski-dense forward orbit?
\subsection{Research attempt}

\paragraph{Scope and source audit (27 September 2026).}
The field is algebraically closed of characteristic zero and may be
countable, including $\overline{\mathbb Q}$. A $K$-point is required, not
merely a generic point over a larger transcendental extension. Dominance
does not make a rational map everywhere defined. In dimension zero the
irreducible variety is a point and the conclusion is immediate; the
discussion below concerns positive dimension.

The supplied Bell--Ghioca paper [S1, Conjecture 1.1] states the general
dichotomy and records the established uncountable-field case, as well as
several special families. Its proposed strengthening for birational maps
of dynamical degree one is not an assumption available in this problem.
The recent Matsuzawa--Xie paper [E1] proves a birational case for projective
varieties over $\overline{\mathbb Q}$ when the first dynamical degree
exceeds the third, including projective threefolds over that field with first
dynamical degree greater than one. Its more general arithmetic-degree
result is not an unconditional dense-orbit theorem for every rational
self-map. Those statements were checked directly in the primary sources;
their deep proofs are not reproduced here.

The present attempt makes the countable-field existence step explicit in
two families: translations on a split torus and simultaneous power maps.
The proofs use finite algebraic identities rather than choosing a point
outside a countable union. The counterexample track tests several tempting
extensions: an arbitrary starting point, a highest possible height-growth
rate, and an assertion that finite avoidance automatically extends to
infinite avoidance. These tests identify failures of methods, not
counterexamples to the catalogue conjecture.

\paragraph{A complete torus-translation criterion.}
Let $T=(\mathbb G_m)^r$, $r\ge1$, and fix
$a=(a_1,\ldots,a_r)\in T(K)$. Consider the everywhere-defined automorphism
\[
 f_a(x_1,\ldots,x_r)=(a_1x_1,\ldots,a_rx_r).
\]
For $u=(u_1,\ldots,u_r)\in\mathbb Z^r$ write
$x^u=\prod x_i^{u_i}$ and $a^u=\prod a_i^{u_i}$. Define
\[
 \Lambda_a=\{u\in\mathbb Z^r:a^u=1\}.
\]
The condition $\Lambda_a=\{0\}$ is multiplicative independence. It also
excludes a nontrivial monomial taking a root-of-unity value: raising that
relation to the order of the root of unity would produce a nonzero
element of $\Lambda_a$.

\textbf{Lemma 484.1.}
For $f_a$ the following conditions are equivalent: $\Lambda_a=0$;
$K(T)^{f_a}=K$; and every point of $T(K)$ has a Zariski-dense forward
orbit. If $\Lambda_a\ne0$, no such orbit is dense.

\emph{Proof of density and the obstruction.}
If $0\ne u\in\Lambda_a$, the nonconstant regular function $x^u$ is
invariant. An orbit starting at $b$ lies in the proper closed hypersurface
$x^u=b^u$, so it is not dense.

Now assume $\Lambda_a=0$. A nonzero Laurent polynomial has the form
\[
 P(x)=\sum_{j=1}^N c_jx^{u_j},
 \qquad c_j\ne0,\quad u_j\text{ distinct}.
\]
Its values on the orbit of $b\in T(K)$ are
\[
 P(f_a^n(b))=\sum_{j=1}^N(c_jb^{u_j})\lambda_j^n,
 \qquad \lambda_j=a^{u_j}.
\]
The $\lambda_j$ are nonzero and pairwise distinct, because an equality
would give $u_i-u_j\in\Lambda_a$. If all orbit values vanished, using
$n=0,\ldots,N-1$ would give a homogeneous system with Vandermonde
determinant
\[
 \prod_{1\le i<j\le N}(\lambda_j-\lambda_i)\ne0.
\]
It would force every $c_jb^{u_j}$ to be zero, a contradiction. Thus no
nonzero Laurent polynomial vanishes on the orbit. Since $T$ is affine
with coordinate ring $K[x_1^{\pm1},\ldots,x_r^{\pm1}]$, this is exactly
Zariski density. The same argument works for any consecutive $N$ iterates,
after multiplying the coefficients by nonzero powers of the $\lambda_j$.

\emph{Proof of the invariant-field assertion.}
Suppose $g=P/Q$ is invariant, with coprime nonzero Laurent polynomials
$P,Q$ in that unique factorization domain. The identity
\[
 P(ax)Q(x)=P(x)Q(ax)
\]
implies $P(x)\mid P(ax)$, by coprimality of $P,Q$. Write
$P(ax)=H(x)P(x)$. For each coordinate, the largest exponent in a product
is the sum of the largest exponents of its factors, and the same is true
of the smallest exponent. This follows by treating the Laurent polynomial
as a one-variable Laurent polynomial over the integral domain generated
by the other coordinates. The exponent ranges of $P(ax)$ and $P(x)$
coincide, so the largest and smallest exponents of $H$ in every coordinate
are both zero. Hence $H$ is a nonzero constant $c$.

Comparing coefficients in $P(ax)=cP(x)$ shows that every exponent $u$ in
the support of $P$ satisfies $a^u=c$. Independence implies that this
support contains only one exponent. Applying the same argument to $Q$
and then the invariance identity gives $Q(ax)=cQ(x)$ with the same $c$.
Thus $P$ and $Q$ are monomials of the same exponent, and $g$ is constant.
If $\Lambda_a\ne0$, the monomial already exhibited is a nonconstant
invariant. This proves all the assertions.\hfill$\square$

For example $a=(2,3,5,\ldots,p_r)$ is independent over every
characteristic-zero field: equality of a rational product of distinct
primes to one forces every exponent to vanish by unique factorization in
$\mathbb Q$. The starting point $(1,\ldots,1)$ therefore works over
$\mathbb Q$ itself, and also over every permitted $K$. Countability causes
no gap in this construction.

\paragraph{A second family with an explicit arithmetic starting point.}
Fix an integer $d\ge2$ and consider
\[
 F_d:T\longrightarrow T,\qquad F_d(x_1,\ldots,x_r)
                              =(x_1^d,\ldots,x_r^d).
\]
Let $q_1,\ldots,q_r$ be distinct positive rational primes and put
$b=(q_1,\ldots,q_r)$. The forward orbit is
$F_d^n(b)=(q_1^{d^n},\ldots,q_r^{d^n})$; it is defined for every $n$.

\textbf{Lemma 484.2.}
This orbit is Zariski dense in $T$ over every algebraically closed
characteristic-zero field $K$. More strongly, a nonzero Laurent polynomial
over $K$ vanishes on only finitely many points of this orbit.

\emph{Proof.}
For $P=\sum_{j=1}^N c_jx^{u_j}$ with nonzero coefficients and distinct
exponents, set $\beta_j=\prod_iq_i^{u_{ji}}$. Unique factorization shows
that the $\beta_j$ are distinct positive rational numbers. Evaluating gives
\[
 P(F_d^n(b))=\sum_{j=1}^N c_j\beta_j^{d^n}.
\]
The coefficients lie in a subfield $k\subset K$ finitely generated over
$\mathbb Q$. Such a field embeds into $\mathbb C$: send a transcendence
basis to algebraically independent complex numbers and extend the map
over the remaining finite algebraic extension. This embedding is used
only for the finitely many coefficients of this one polynomial; there is
no assertion that the entire field $K$ embeds into $\mathbb C$.

After applying the embedding, order the rational bases as
$0<\beta_1<\cdots<\beta_N$. Division by $\beta_N^{d^n}$ gives
\[
 \beta_N^{-d^n}P(F_d^n(b))
 =c_N+\sum_{j<N}c_j(\beta_j/\beta_N)^{d^n}\longrightarrow c_N\ne0.
\]
Consequently the expression is nonzero for all sufficiently large $n$.
Injectivity of the embedding transfers this conclusion back to $K$.
Every proper closed subset of the affine torus is annihilated by a nonzero
Laurent polynomial, so the orbit meets it only finitely often. In
particular the orbit and all its tails are dense.\hfill$\square$

This also verifies the invariant-field hypothesis in this family. If
$g=P/Q\in K(T)$ satisfies $g\circ F_d=g$, the denominator $Q$ is nonzero
on all sufficiently late orbit points by Lemma 484.2. On that tail the
rational identity can be evaluated at each consecutive pair, so $g$ has a
constant value $c$. Density of the tail implies $P-cQ=0$ as a Laurent
polynomial, and $g=c$. This argument checks denominator avoidance rather
than silently evaluating a rational function at its base points.

The proof is insensitive to whether the largest base is below or above
one: it is the ratios $\beta_j/\beta_N<1$ that matter. Negative Laurent
exponents therefore cause no problem. In dimension one it reduces to the
simple fact that the sequence $q^{d^n}$ is infinite and hence dense in
$\mathbb G_m$.

\paragraph{A maximal-growth orbit still need not be dense.}
Even in the family just solved, the hypothesis does not make every
starting point suitable. For $F_2$ on $(\mathbb G_m)^2$, the orbit of
$(2,2)$ is contained in the diagonal $x_1=x_2$, whereas the orbit of
$(2,3)$ is dense by Lemma 484.2. Both have the same exponential growth
rate for logarithmic projective height. With
$h(x_1,x_2)=h_{\mathbb P^1}(x_1)+h_{\mathbb P^1}(x_2)$ on rational
positive coordinates,
\[
 h(F_2^n(2,2))=2^{n+1}\log2,
 \qquad \lim_{n\to\infty}\max\{1,h(F_2^n(2,2))\}^{1/n}=2.
\]
On $\mathbb P^1\times\mathbb P^1$ the map pulls back each ruling
divisor to twice itself, so its first dynamical degree is also two.
Thus even attaining the largest predicted height-growth rate does not,
by itself, imply a dense orbit. The additional geometry in the recent
results cited above cannot be replaced by that implication.

Roots of unity give a second useful stress test. On $\mathbb G_m$ the
map $x\mapsto x^2$ has trivial invariant rational-function field, but
every root of unity has a finite forward orbit. There are infinitely many
such roots, and an infinite subset of an irreducible curve is Zariski
dense. Consequently no nonempty Zariski-open subset consists entirely of
dense-orbit points in this example. The conjecture asks for existence of
one point; it does not assert that deleting a single algebraic exceptional
set leaves only successful starting points.

\paragraph{Finite avoidance is valid; countable avoidance is an extra step.}
Let $f:X\dashrightarrow X$ be an arbitrary dominant rational self-map.
For each integer $N\ge0$, let $U_N$ be the set of points for which the
first $N$ applications of $f$ are successively defined. This is a dense
open subset. Indeed $U_0=X$ and the induction step restricts a dominant
iterate to the inverse image of the nonempty open domain of $f$; that
inverse image is nonempty open because the iterate is dominant. A finite
intersection of these domains is still nonempty by irreducibility.

Since $K$ is algebraically closed, each nonempty open $U_N$ has a
$K$-point. This proves
\[
 \text{for every }N\text{ there exists }x_N\in X(K)
 \text{ whose orbit is defined through time }N.
\]
It does not prove existence of a single $x$ that works for all $N$.
The required order of quantifiers is reversed, and there is no compactness
of the set of $K$-points in the Zariski topology that repairs that reversal.

For a model of the failure over a countable field, enumerate
$\overline{\mathbb Q}=\{a_1,a_2,\ldots\}$. The subsets
\[
 V_N=\mathbb A^1(\overline{\mathbb Q})\setminus\{a_1,\ldots,a_N\}
\]
are nonempty open and decreasing, but their intersection has no
$\overline{\mathbb Q}$-point. These particular sets are not claimed to be
the orbit domains of a rational self-map. They disprove the proposed
general countable-avoidance principle, which would otherwise be used
without any dynamical justification.

Similarly, over a countable algebraically closed field the union of all
closed points is a countable union of proper closed subsets covering
$\mathbb A^1(K)$. Passing to an uncountable extension and choosing a
point outside that union does not descend the point to $K$. All varieties
and maps can have finitely many defining coefficients while the property
of avoiding every iterate or every proper orbit closure remains an
infinite collection of conditions on the starting point.

\paragraph{Why a finite orbit computation does not close the argument.}
Noetherianity says descending chains of closed subsets stabilize; it does
not say that ascending chains of finite orbit prefixes stabilize. The
closed subsets $\{1,\ldots,N\}\subset\mathbb A^1$ are an explicit
ascending nonstationary chain. Thus taking larger finite prefixes,
computing their algebraic relations, and appealing to Noetherianity cannot
certify the final orbit closure without an independent degree bound or
recurrence argument.

The Vandermonde proof succeeds for torus translations because every
Laurent polynomial with $N$ terms is ruled out by $N$ consecutive iterates,
and the argument works symbolically for arbitrary $N$. The power-map
proof instead uses an asymptotic separation of finitely many distinct
positive bases, separately for every polynomial. Neither proof gives
these properties for an arbitrary rational map. The lack of invariant
rational functions has not been converted here into an equally concrete
mechanism excluding every lower-dimensional orbit closure.

One possible next statement would be a dynamical avoidance lemma: from
$K(X)^f=K$, construct a $K$-point with an infinite defined orbit and prove
that every proper closed set misses at least one point of that orbit.
The substantive missing input is a means of selecting that point over the
original countable field while controlling all iterates. It cannot be
replaced by the finite nonempty-open argument just given. An alternative
would provide effective algebraic restrictions on all possible proper
orbit closures, strong enough that finitely many explicitly verified
conditions exclude them. No such restriction is established in general.

\paragraph{Exact checks and outcome.}
The companion script computes an exact Vandermonde certificate for the
translation $(x,y)\mapsto(2x,3y)$ on the 16 monomials $x^iy^j$ with
$0\le i,j\le3$. It compares the product-of-differences determinant with
an independent fraction-free elimination. A second exact evaluation
matrix tests the first nine iterates of $(2,3)$ under the squaring map
against the nine monomials with $0\le i,j\le2$. Nonzero determinants
exclude relations in those finite spaces; they are not a computational
proof of density by themselves. The all-polynomial density assertions
are proved above. All starting points and coefficients in the finite
certificates are rational, so no generic transcendental point is hidden.

\textbf{Outcome: UNRESOLVED, with complete elementary families.}
The attempt proves the torus-translation criterion, constructs explicit
dense power-map orbits over the permitted countable fields, and identifies
specific failures of arbitrary starting-point selection, growth-only tests,
and finite-to-infinite avoidance. The first remaining gap is a construction
of one original-field point whose defined orbit excludes every proper
closed closure for an arbitrary dominant rational map with trivial
invariant field. No general proof or counterexample is supplied.

\subsection{Sources}
\begin{itemize}
\item[S1]Jason Bell and Dragos Ghioca, A conjecture strengthening the Zariski dense orbit problem for birational maps of dynamical degree one.
\url{https://personal.math.ubc.ca/~dghioca/papers/conjecture_automorphisms.pdf}.
\textit{Formulation source.}
\item[E1]Yohsuke Matsuzawa and Junyi Xie, \emph{Arithmetic degree
and its application to Zariski dense orbit conjecture}, Journal of the
London Mathematical Society 112 (2025), e70282;
author manuscript arXiv:2409.06160v2, 17 September 2024.
\url{https://doi.org/10.1112/jlms.70282};
\url{https://arxiv.org/html/2409.06160v2}.
\textit{Publisher abstract and manuscript introduction inspected
27 September 2026. Birational and dynamical-degree hypotheses retained;
the height-growth theorem is not used as a general dense-orbit theorem.}
\end{itemize}
\end{document}
